Os resultados observados no estado atual do repositório devem ser lidos como consequência direta dos capítulos anteriores. Depois de estabelecer os fundamentos matemáticos e a separação arquitetural entre cadeia e componentes externos, torna-se possível avaliar se o sistema preserva coerência entre teoria, implementação e operação. Nesse sentido, o primeiro resultado é funcional: o sistema é capaz de codificar arquivos, gerar metadados, distribuir pacotes e reconstruir a carga útil original a partir das estruturas produzidas na codificação. Esse comportamento confirma que o encadeamento entre divisão em blocos, redundância por XOR e resolução iterativa de pacotes é operacionalmente consistente com a proposta do HySail~\cite{hysail_readme}.

Em segundo lugar, a implementação mostra que a auditoria baseada em resto de divisão polinomial foi incorporada ao fluxo real do sistema. A função \texttt{gf2\_poly\_mod} não permanece apenas como construção teórica; ela participa efetivamente da geração de MACs e, por consequência, da lógica que sustenta o mecanismo de desafio-resposta. Os testes do projeto reforçam essa observação ao validar largura do resto, tratamento de divisor nulo e consistência de resultados conhecidos. Isso confirma a conexão defendida no capítulo matemático: a representação em $GF(2)$ não é decorativa, mas operacional.

Em terceiro lugar, a integração com blockchain produz um resultado arquitetural relevante: o repositório demonstra que é viável separar reconstrução e liquidação. A reconstrução ocorre fora da cadeia, enquanto os contratos registram arquivos, provedores, solicitações de download, aceitação de blocos e reembolso de saldos. Essa divisão é coerente com as restrições de Ethereum e sugere que a auditoria do HySail pode ser complementada por uma camada descentralizada de coordenação sem que a cadeia assuma o custo de manipular o arquivo completo~\cite{hysail_dapp_architecture}. Em outras palavras, o fluxo mermaid adotado em doc não é apenas um diagrama conceitual; ele corresponde a uma decomposição que faz sentido técnico para a implementação atual.

Por fim, o principal resultado analítico desta tese é a identificação de uma diferença entre a formulação original e a implementação atual. A base conceitual do Hy-SAIL permanece presente, sobretudo na relação entre disponibilidade, corpos finitos e autenticação homomórfica. Entretanto, a codificação efetivamente implementada aproxima-se mais de um esquema inspirado em LT Code do que de uma realização integral de Online Code com blocos auxiliares explícitos. Esse resultado não invalida a proposta; ao contrário, esclarece o ponto exato em que a implementação atual simplifica a formulação original e, com isso, delimita um caminho objetivo para evoluções futuras.
